% !TEX program = pdflatex
% =====================================================================
% Zwätschgen-Pubquiz -- Konzept
% Der Fragenkatalog liegt separat in fragenkatalog.tex
% =====================================================================
\documentclass[a4paper,11pt]{article}
% =====================================================================
% Gemeinsame Präambel für main.tex (Konzept) und fragenkatalog.tex
% Einbinden mit % =====================================================================
% Gemeinsame Präambel für main.tex (Konzept) und fragenkatalog.tex
% Einbinden mit % =====================================================================
% Gemeinsame Präambel für main.tex (Konzept) und fragenkatalog.tex
% Einbinden mit \input{praeambel}
% =====================================================================
\usepackage[utf8]{inputenc}
\usepackage[T1]{fontenc}
\usepackage{lmodern}
\usepackage[ngerman]{babel}
\usepackage[margin=2.3cm]{geometry}
\usepackage{enumitem}
\usepackage{booktabs}
\usepackage{tabularx}
\usepackage{xcolor}
\usepackage{listings}
\usepackage{fancyhdr}
\usepackage{amsmath}
\usepackage{array}
\usepackage[hidelinks]{hyperref}

% ---------- Farben ----------
\definecolor{zwetschge}{HTML}{5B2A86}
\definecolor{offen}{HTML}{C0392B}

% ---------- Kopf- und Fusszeile ----------
\pagestyle{fancy}
\fancyhf{}
\lhead{Zwätschgenabend -- Pubquiz}
\rhead{HSLU Elektrotechnik}
\cfoot{\thepage}
\setlength{\headheight}{14pt}

% ---------- Listen ----------
\setlist[itemize]{itemsep=0.2em, topsep=0.3em}

% ---------- Hilfsmakros ----------
% Markiert Punkte, die das OK noch klären muss
\newcommand{\offen}[1]{\textcolor{offen}{\textbf{[offen:} #1\textbf{]}}}

\lstset{
  basicstyle=\ttfamily\small,
  frame=single,
  rulecolor=\color{zwetschge!50},
  language=C,
  columns=fullflexible,
  xleftmargin=1em
}

% =====================================================================
\usepackage[utf8]{inputenc}
\usepackage[T1]{fontenc}
\usepackage{lmodern}
\usepackage[ngerman]{babel}
\usepackage[margin=2.3cm]{geometry}
\usepackage{enumitem}
\usepackage{booktabs}
\usepackage{tabularx}
\usepackage{xcolor}
\usepackage{listings}
\usepackage{fancyhdr}
\usepackage{amsmath}
\usepackage{array}
\usepackage[hidelinks]{hyperref}

% ---------- Farben ----------
\definecolor{zwetschge}{HTML}{5B2A86}
\definecolor{offen}{HTML}{C0392B}

% ---------- Kopf- und Fusszeile ----------
\pagestyle{fancy}
\fancyhf{}
\lhead{Zwätschgenabend -- Pubquiz}
\rhead{HSLU Elektrotechnik}
\cfoot{\thepage}
\setlength{\headheight}{14pt}

% ---------- Listen ----------
\setlist[itemize]{itemsep=0.2em, topsep=0.3em}

% ---------- Hilfsmakros ----------
% Markiert Punkte, die das OK noch klären muss
\newcommand{\offen}[1]{\textcolor{offen}{\textbf{[offen:} #1\textbf{]}}}

\lstset{
  basicstyle=\ttfamily\small,
  frame=single,
  rulecolor=\color{zwetschge!50},
  language=C,
  columns=fullflexible,
  xleftmargin=1em
}

% =====================================================================
\usepackage[utf8]{inputenc}
\usepackage[T1]{fontenc}
\usepackage{lmodern}
\usepackage[ngerman]{babel}
\usepackage[margin=2.3cm]{geometry}
\usepackage{enumitem}
\usepackage{booktabs}
\usepackage{tabularx}
\usepackage{xcolor}
\usepackage{listings}
\usepackage{fancyhdr}
\usepackage{amsmath}
\usepackage{array}
\usepackage[hidelinks]{hyperref}

% ---------- Farben ----------
\definecolor{zwetschge}{HTML}{5B2A86}
\definecolor{offen}{HTML}{C0392B}

% ---------- Kopf- und Fusszeile ----------
\pagestyle{fancy}
\fancyhf{}
\lhead{Zwätschgenabend -- Pubquiz}
\rhead{HSLU Elektrotechnik}
\cfoot{\thepage}
\setlength{\headheight}{14pt}

% ---------- Listen ----------
\setlist[itemize]{itemsep=0.2em, topsep=0.3em}

% ---------- Hilfsmakros ----------
% Markiert Punkte, die das OK noch klären muss
\newcommand{\offen}[1]{\textcolor{offen}{\textbf{[offen:} #1\textbf{]}}}

\lstset{
  basicstyle=\ttfamily\small,
  frame=single,
  rulecolor=\color{zwetschge!50},
  language=C,
  columns=fullflexible,
  xleftmargin=1em
}


\title{\textcolor{zwetschge}{\textbf{Zwätschgen-Pubquiz}}\\[0.3em]\large Konzept}
\author{OK Zwätschgenabend, HSLU Elektrotechnik}
\date{\today}

\begin{document}
\maketitle
\thispagestyle{fancy}
\tableofcontents

% ---------------------------------------------------------------------
\section{Ziel und Grundidee}
Am Zwätschgenabend kommen alle ETler der HSLU zusammen. Dieses Jahr gibt es zusätzlich ein Pubquiz. Es verfolgt drei Ziele:
\begin{itemize}
  \item \textbf{Jahrgänge mischen:} Die Teams werden so gebildet, dass Studierende aus allen drei Jahrgängen an einem Tisch sitzen.
  \item \textbf{Ins Gespräch kommen:} Die Fragen sind so schwer, dass ein Jahrgang allein sie nicht lösen kann. Das Team muss sich austauschen.
  \item \textbf{Dozierende einbinden:} Dozierende spielen mit, ohne dass sich Studierende und Dozierende in getrennten Lagern gegenübersitzen.
\end{itemize}

% ---------------------------------------------------------------------
\section{Rahmenbedingungen}
\begin{tabularx}{\textwidth}{@{}l X@{}}
\toprule
Datum & 3.12.2026\\
Ort & IHomeLab \\
Teilnehmende & ca. 60 \\
Dauer Quiz & ca.\ 1,5 Stunden inklusive Pause \\
Technik & Beamer, Lautsprecher, Mikrofon \\
\bottomrule
\end{tabularx}

% ---------------------------------------------------------------------
\section{Teambildung}
\subsection{Jahrgangsmix}
\begin{itemize}
  \item Beim Eintreffen zieht jede Person eine nummerierte Zwetschgenkarte. Gleiche Nummer heisst gleiches Team.
  \item Die Karten werden vorab pro Jahrgang vorsortiert (z.\,B.\ drei Stapel in verschiedenen Farben). So sitzen pro Team höchstens zwei Personen aus demselben Jahrgang.
  \item Zielgrösse: 5 bis 6 Personen pro Team, ideal 2--2--2.
\end{itemize}

\subsection{Teamgrösse nach Teilnehmerzahl}
\begin{tabularx}{\textwidth}{@{}l l X@{}}
\toprule
\textbf{Teilnehmende} & \textbf{Teams} & \textbf{Bemerkung} \\
\midrule
30 -- 40 & 6 -- 7 & Auswertung von Hand gut machbar \\
40 -- 60 & 8 -- 10 & Zweite Person für Auswertung einplanen \\
über 60 & 10 -- 12 & Antwortblätter tischweise einsammeln \\
\bottomrule
\end{tabularx}


% ---------------------------------------------------------------------
\section{Einbindung der Dozierenden}
Die Dozierenden sollen dabei sein, ohne dass sie Zusatzaufwand haben. Beide Varianten kommen ohne Rotieren oder Tischwechsel aus: hinsetzen, mitraten, trinken.

\medskip
\noindent\begin{tabularx}{\textwidth}{@{}>{\raggedright\arraybackslash}p{2.2cm} >{\raggedright\arraybackslash}X >{\raggedright\arraybackslash}X@{}}
\toprule
 & \textbf{Variante A: Dozententeam} & \textbf{Variante B: Dozierende im Studententeam} \\
\midrule
Prinzip & Die Dozierenden bilden ein eigenes Team und spielen regulär mit. & Jedes Studententeam bekommt beim Losen einen Dozierenden fest zugeteilt. Er sitzt den ganzen Abend am selben Tisch. \\[0.3em]
Wertung & Das Dozententeam steht normal in der Rangliste. Jedes Studententeam, das am Ende vor den Dozierenden liegt, erhält einen Bonus (z.\,B.\ eine Runde Getränke). & Normale Wertung. Die Dozierenden zählen als gewöhnliches Teammitglied. \\[0.3em]
Vorteil & Klarer «Gegner», den alle schlagen wollen. Die Dozierenden sitzen unter sich, was vielen angenehm ist. & Studierende und Dozierende kommen ins Gespräch, die Teams bleiben gemischt. \\[0.3em]
Nachteil & Während des Quiz kein Austausch zwischen Dozierenden und Studierenden. & Braucht etwa so viele Dozierende wie Teams (bei 60 Personen ca.\ 10). Teams mit dem passenden Fachdozierenden haben in einzelnen Runden einen Vorteil. \\
\bottomrule
\end{tabularx}

\medskip
\noindent\textbf{Empfehlung:} Die Wahl hängt davon ab, wie viele Dozierende zusagen. Kommen etwa so viele wie es Teams gibt, ist Variante B die bessere Wahl, weil sie den Austausch fördert. Sind es deutlich weniger, wird Variante A gespielt. Variante A ist auch die Rückfallebene, falls kurzfristig Dozierende absagen.

% ---------------------------------------------------------------------
\section{Ablauf grob angedacht}
\begin{tabularx}{\textwidth}{@{}l l X@{}}
\toprule
\textbf{Zeit} & \textbf{Block} & \textbf{Inhalt} \\
\midrule
0:00 & Start & Begrüssung, Regeln, Teams finden \\
0:05 & Runde 1 & Elektrotechnik mit Tücken \\
0:15 & Runde 2 & Digitaltechnik und Embedded \\
0:25 & Runde 3 & Signale, Regelung, Kommunikation \\
0:35 & Pause & Essen und Trinken, Zwischenstand \\
0:45 & Runde 4 & Energie und Schweiz \\
0:53 & Runde 5 & Geschichte und Köpfe \\
1:01 & Runde 6 & Schätzfragen \\
1:09 & Runde 7 & Zwätschgen-Spezial \\
1:15 & Finale & Finalfrage mit Punkteinsatz \\
1:25 & Ende & Siegerehrung \\
\bottomrule
\end{tabularx}

% ---------------------------------------------------------------------
\section{Rundenformate}
Der vollständige Fragenkatalog liegt in \texttt{fragenkatalog.tex}. Hier sind nur die Formate beschrieben.

\begin{description}[style=nextline, leftmargin=1.5em]
  \item[Multiple-Choice-Runden (Runden 1--5, 7)] Pro Runde 8 bis 10 Fragen mit je vier Antwortmöglichkeiten A--D. Die Fragen sind anspruchsvoll, lassen sich aber mit Überlegen, Rechnen oder Ausschlussverfahren lösen. Die falschen Antworten sind bewusst plausibel gewählt.
  \item[Schätzfragen (Runde 6)] Offene Fragen mit einer Zahl als Antwort. Gewertet wird die Nähe zur richtigen Lösung.
  \item[Bild- und Hörrunde (optional)] Schaltzeichen, Platinen, Geräusche (Modem, Netzbrummen, Relais). Diese Runde kann bei Zeitmangel entfallen.
  \item[Dozentenrunde (optional)] Kinderfotos zuordnen, Vorlesungszitate erraten, «Zwei Wahrheiten, eine Lüge». Die Dozierenden müssen vorab Material liefern.
  \item[Teamaktion] Widerstandsfarbcode-Speedrun oder Spaghetti-Turm mit Zeitlimit.
  \item[Finalfrage] Eine Frage, bei der jedes Team vorab Punkte setzt.
\end{description}

% ---------------------------------------------------------------------
\section{Wertung}
\begin{itemize}
  \item \textbf{Multiple Choice:} 2 Punkte pro richtiger Antwort, keine Minuspunkte.
  \item \textbf{Schätzfragen:} Das nächstliegende Team erhält 3 Punkte, das zweitnächste 2, das drittnächste 1. Bei Werten über mehrere Grössenordnungen wird der Abstand logarithmisch bestimmt (Faktor statt Differenz).
  \item \textbf{Teamaktion:} 5, 3 und 1 Punkte für die Plätze 1 bis 3.
  \item \textbf{Gleichstand:} Eine zusätzliche Schätzfrage entscheidet.
  \item Handys bleiben während der Runden in der Tasche.
\end{itemize}


% ---------------------------------------------------------------------
\section{Material und Technik}
\begin{itemize}
  \item Zwetschgenkarten für die Teameinteilung (pro Jahrgang eine Farbe)
  \item Antwortblätter pro Runde und Team (Vorlage im Fragenkatalog), Stifte
  \item Beamer mit Folien zu den Fragen, Lautsprecher für die Hörrunde
  \item Live-Punktestand als Google Sheet am Beamer
  \item Material für die Teamaktion (Widerstände, Spaghetti, Klebeband)
  \item Preise für die Plätze 1 bis 3, Trostpreis und Preis für den besten Teamnamen
\end{itemize}

% ---------------------------------------------------------------------
\section{Vorbereitung}
\begin{tabularx}{\textwidth}{@{}l X@{}}
\toprule
\textbf{Wer} & \textbf{Was} \\
\midrule
Tom & Datum und Raum fixieren, Dozierende anfragen \\
Chosinus & Fragenkatalog finalisieren, zeitabhängige Fakten prüfen \\
Chosinus & Folien erstellen, Preise organisieren \\
Chosinzs & Material vorbereiten, Technik testen, Antwortblätter verteilen \\
Chosinus & Punktestand live führen, Siegerehrung vorbereiten. Preise? \\
\bottomrule
\end{tabularx}

% ---------------------------------------------------------------------
\section{Preise }
\begin{itemize}
  \item 1.\ Platz: Zwetschgenschnaps
  \item 2.\ und 3.\ Platz: HotDog oder ka
  \item Letzter Platz: ein Sack Widerstände
\end{itemize}


\end{document}
